\section{Related work}
\label{sec:related}

Probes establish that a variable can be recovered from a representation, but not that the model's
dynamics use or preserve it \citep{alain2017probing}. Control tasks \citep{hewitt2019control} and
later critiques \citep{belinkov2022probing} show that probe accuracy mixes what the representation
stores with what the probe can learn, and our untrained baseline is the same problem in a different
form: four of six randomly initialized models exceed the threshold we had previously used. Work on
emergent world representations goes further by intervening on the representation
\citep{li2023othello,nanda2023othello}, which is the logic our rollout correction follows.

A separate literature imposes physical structure by design. Equation-discovery methods such as SINDy
\citep{brunton2016sindy} infer governing equations from observed trajectories, and Hamiltonian and
Lagrangian networks \citep{greydanus2019hnn,cranmer2020lnn} build conservation into an architecture
so that it holds by construction. We instead ask whether such structure emerges in a conventionally
trained world model, and analyze it after the fact using standard linear-algebraic tools related to
Koopman and dynamic mode methods \citep{schmid2010dmd}.

World models are normally evaluated through prediction quality or control performance
\citep{hafner2023dreamerv3,ha2018world,schrittwieser2020muzero}. Those metrics do not reveal whether
the latent transition preserves a physical constraint it has learned. Our intervention tests whether
violating a recovered latent constraint itself contributes to rollout error, and the damped arm
tests whether the extraction procedure returns a conserved-looking scalar when the training dynamics
have no such invariant. The disentanglement literature \citep{locatello2019challenging} documents how
readily an appealing structural claim about representations survives without that kind of control.
